\section{Problem and basic properties}\label{sec:foundations}

A relaxed switching control allocates one unit of activity among finitely many
modes. Rounding selects one mode at a time. The cumulative allocation error is
the objective of combinatorial integral approximation (CIA), the rounding
subproblem in a decomposition of mixed-integer optimal control. A restriction
on the number of mode changes is a hard constraint on that subproblem.
Sager and Zeile~\cite{sager-zeile2021} study this formulation, develop maximum
dwell rounding, and give sharp bounds and a conjecture for its worst-case
error. Our focus is the extremal rounding problem itself. In particular, a
bound on cumulative allocation error is not an assertion that the same
schedule optimizes an underlying nonlinear control objective.

\subsection{Controls, schedules, and the two error criteria}

Fix a horizon $T>0$ and an integer number of modes $n\ge2$. Throughout,
switch budgets $s$ are nonnegative integers, activation-block counts $k$ are
positive integers, and grid sizes $N$ are positive integers. Write
$[n]=\{1,\ldots,n\}$ and
\[
 \Delta_n=\left\{a\in\R^n:a_i\ge0,\ \sum_{i=1}^n a_i=1\right\}.
\]
A relaxed control is a Lebesgue-measurable map
$\alpha:[0,T]\to\Delta_n$, considered up to equality almost everywhere.
Its cumulative allocation and terminal masses are
\begin{equation}\label{eq:cumulatives}
 A_i(t)=\int_0^t\alpha_i(u)\,du,\qquad m_i=A_i(T).
\end{equation}
Thus $A_i$ is nondecreasing and $1$-Lipschitz, $A_i(0)=0$, and
$\sum_i A_i(t)=t$. Conversely, every vector of functions with these properties
is the cumulative allocation of such a control: Lipschitz functions are
absolutely continuous and differentiable almost everywhere, their derivatives
are nonnegative, and differentiation of the sum gives a simplex-valued
vector almost everywhere. Denote this class of cumulative allocations by
$\Acal_n(T)$. We use either $\alpha$ or $A$ to specify an input.

An integer control, or schedule, selects a coordinate vector $e_i$ on each of
finitely many intervals. Its initial activation is free: only changes between
two successive active modes count as switches. Its values at isolated times
are immaterial. Let $\Wcal_{n,k}(T)$ be the set of cumulative occupations
\[
 W_i(t)=\int_0^t\omega_i(u)\,du
\]
of schedules with at most $k$ positive-length activation blocks, equivalently
at most $s=k-1$ switches. Unless explicitly stated otherwise, repeated modes
are allowed, the initial and final modes are unrestricted, and there are no
dwell-time or transition restrictions. A parameterization of this entire set
uses a word $p\in[n]^k$ and ordered times
\begin{equation}\label{eq:schedule-parameterization}
 0=t_0\le t_1\le\cdots\le t_k=T,\qquad
 \omega(t)=e_{p_j}\quad(t_{j-1}<t<t_j).
\end{equation}
Zero-length intervals can be deleted and consecutive equal modes can be
merged. These operations cannot increase the switch count. Consequently this
parameterization includes shorter schedules without imposing distinctness
on competing modes.

The full error and the one-sided error are respectively
\begin{equation}\label{eq:errors}
 D(A,W)=\max_{i\in[n]}\max_{0\le t\le T}|A_i(t)-W_i(t)|,
 \qquad
 D^-(A,W)=\max_{i\in[n]}\max_{0\le t\le T}(W_i(t)-A_i(t)).
\end{equation}
The minus superscript refers to the negative part of the discrepancy $A-W$;
the displayed one-sided maximum is nonnegative because it includes $t=0$.
For a fixed input, define the continuous instance optima
\begin{equation}\label{eq:instance-optima}
 \OPT_s(A)=\min_{W\in\Wcal_{n,s+1}(T)}D(A,W),\qquad
 \OPT^-_k(A)=\min_{W\in\Wcal_{n,k}(T)}D^-(A,W).
\end{equation}
The corresponding minimax values are
\begin{equation}\label{eq:minimax-values}
 F_{n,s}(T)=\sup_{A\in\Acal_n(T)}\OPT_s(A),\qquad
 G^-_{n,k}(T)=\sup_{A\in\Acal_n(T)}\OPT^-_k(A).
\end{equation}
The indexing is intentional: $F$ counts switches, whereas $G^-$ counts
activation blocks, with $k=s+1$. Both infima in~\eqref{eq:instance-optima}
and both suprema in~\eqref{eq:minimax-values} are attained, as shown below.

\subsection{Fixed grids and endpoint evaluation}

A grid is a strictly increasing sequence
$\mathcal T=(x_0,\ldots,x_N)$ with $x_0=0$, $x_N=T$. Its cell lengths are
$\Delta_j=x_j-x_{j-1}$ and $\bar\Delta=\max_j\Delta_j$.
Let $\Acal_n(\mathcal T)$ consist of cumulative allocations of controls
constant on each cell, and let $\Wcal_{n,k}(\mathcal T)$ restrict the switching
times to grid points. Set
\begin{equation}\label{eq:grid-definitions}
 \OPT_s^{\mathcal T}(A)=
 \min_{W\in\Wcal_{n,s+1}(\mathcal T)}D(A,W),\qquad
 F^{\mathcal T}_{n,s}=
 \max_{A\in\Acal_n(\mathcal T)}\OPT_s^{\mathcal T}(A).
\end{equation}
The instance optimum is defined even for $A\in\Acal_n(T)$ that is not
piecewise affine. The grid minimax in~\eqref{eq:grid-definitions} does restrict
the input to cellwise constant controls. Analogously, use
$\OPT_k^{-,\mathcal T}$ and $G_{n,k}^{-,\mathcal T}$ for the one-sided
criterion. On the unit grid $(0,1,\ldots,N)$ we abbreviate the full minimax by
$F^{\mathrm{unit}}_{n,s}(N)$. Every grid with cell length $\Delta$ is obtained
from a unit grid by scaling all errors by $\Delta$.

For the same input, $\OPT_s(A)\le\OPT_s^{\mathcal T}(A)$. This fact alone
does not compare the two full minimax values, whose input classes differ.
The following observation permits exact evaluation even for merely measurable
inputs and will be used throughout.

\begin{lemma}[Endpoint monotonicity]\label{lem:endpoint}
On any interval on which the selected mode is $p$, the discrepancy
$A_p-W_p$ is nonincreasing and every other discrepancy $A_i-W_i$ is
nondecreasing. Hence both errors in~\eqref{eq:errors} are determined by the
activation-block endpoints. In particular, for a grid schedule they are
determined by grid endpoints, even when $\alpha$ varies inside grid cells.
\end{lemma}
\begin{proof}
The absolutely continuous discrepancies have derivatives $\alpha_p-1\le0$
and $\alpha_i\ge0$ for $i\ne p$, almost everywhere on the interval. Integrating
these inequalities proves monotonicity. A monotone real function takes its
largest absolute value and each of its one-sided extrema at an endpoint.
\end{proof}

Replacing a measurable control by its cell averages preserves every grid
endpoint $A_i(x_j)$. Lemma~\ref{lem:endpoint} therefore shows that this
replacement preserves the error of every grid schedule, and hence its grid
instance optimum. This observation does not say that averaging preserves
its continuous instance optimum. It does, however, give the minimax comparison
\begin{equation}\label{eq:minimax-grid-domination}
 F_{n,s}(T)\le F^{\mathcal T}_{n,s},\qquad
 G^-_{n,k}(T)\le G^{-,\mathcal T}_{n,k}.
\end{equation}
Indeed, for any measurable input $A$, its cell average $\bar A$ satisfies
$\OPT_s(A)\le\OPT_s^{\mathcal T}(A)=
\OPT_s^{\mathcal T}(\bar A)\le F^{\mathcal T}_{n,s}$.
Taking the supremum proves the first inequality; the one-sided argument is
identical.

\subsection{Existence, continuity, and scaling}

Equip vector-valued cumulative functions with
$\norm{A}_\infty=\max_i\max_{t\in[0,T]}|A_i(t)|$.
\begin{proposition}[Attainment and stability]\label{prop:compactness}
The sets $\Acal_n(T)$ and $\Wcal_{n,k}(T)$ are compact in the uniform norm.
Each of the four instance optimum maps defined above is $1$-Lipschitz in
$A$. All continuous and grid instance optima and minimax values are attained.
\end{proposition}
\begin{proof}
The cumulative allocations are uniformly bounded and equicontinuous. Their
uniform limits preserve monotonicity, the Lipschitz constants, initial values,
and the sum identity, and therefore belong to $\Acal_n(T)$ by the converse
following~\eqref{eq:cumulatives}. The Arzel\`a--Ascoli theorem gives compactness.

For a fixed word $p$, the ordered time simplex in
\eqref{eq:schedule-parameterization} is compact. Its cumulative occupations
are continuous in those times. Indeed, the useful representation
\[
 W_i(t)=\mathbf1_{\{p_k=i\}}t+
 \sum_{j=1}^{k-1}
 (\mathbf1_{\{p_j=i\}}-\mathbf1_{\{p_{j+1}=i\}})\min\{t,t_j\}
\]
gives
$\norm{W(\cdot;t_1,\ldots,t_{k-1})-
W(\cdot;t'_1,\ldots,t'_{k-1})}_\infty
\le\sum_{j=1}^{k-1}|t_j-t'_j|$.
Taking the finite union over words proves compactness of
$\Wcal_{n,k}(T)$. Both error criteria are continuous, so their minima are
attained. For either criterion $d$, the inequality
$|d(A,W)-d(B,W)|\le\norm{A-B}_\infty$ holds for every $W$. Taking minima
shows the stated Lipschitz property. Compactness of $\Acal_n(T)$ then gives
attainment of the outer maxima. For grids, there are finitely many schedules
and the input is a product of $N$ compact simplexes, proving the analogous
claims.
\end{proof}

Changing time by $t=Tu$ is a bijection of both input and schedule classes and
multiplies errors by $T$. Thus
\begin{equation}\label{eq:scaling}
 F_{n,s}(T)=T F_{n,s}(1),\qquad G^-_{n,k}(T)=T G^-_{n,k}(1).
\end{equation}
The same statement applies to instance optima and a simultaneously rescaled
grid. At $T=0$ all errors are zero; results stated for $T>0$ extend with this
convention. Enlarging the switch budget cannot increase any instance optimum
or minimax value.

\begin{proposition}[No-switch baseline]\label{prop:no-switch}
For every $n\ge2$,
$F_{n,0}(T)=T(1-1/n)$. On any grid,
$F^{\mathcal T}_{n,0}=T(1-1/n)$ as well.
\end{proposition}
\begin{proof}
Activating mode $p$ throughout gives error
$\max\{T-m_p,\max_{i\ne p}m_i\}=T-m_p$, since
$\sum_{i\ne p}m_i=T-m_p$. Choosing a largest terminal mass gives error at
most $T(1-1/n)$. The uniform control $\alpha_i=1/n$ gives equality and
belongs to every grid input class.
\end{proof}
